\section{第14版の計算証明書と再現範囲}
\label{app:v14-reproducibility}

第14版の追加ファイルは継続用bundleの \texttt{source\_v14/} に置き、
添付由来の69ファイルを含む \texttt{source\_v13/} は変更していない。
新規ファイルは \texttt{characters/}, \texttt{rank/}, \texttt{contact/},
\texttt{detector/}, \texttt{tate/} に分けた。
日本語原稿はLuaLaTeXで組版する。

指標計算は、Euler積による整数の重み計算をNewton--Adams再帰と照合した。
通常接触階数は、一般の整数三次テンソルでの94方向と60方向の独立性に基づく。
両者は別の辺番号付け・縮約方式でも監査し、選択minorの全成分を有理数で再計算した。
五点汎関数は300個の通常接触値と80個の追加方向の値を新規に縮約して作った。
これらの零判定は有理数で行っている。

実際の18次係数は、公開した五点それぞれからTate再帰を初めから作った。
標準点では通常のFraction演算と共通分母整数演算を照合した。
全ての点で、閉性と切断されたMaurer--Cartan等式、
独立な定数ブロック抽出、二方式の行列式係数計算を確認した。
この新規再構成と、保存した行列・値を読む高速監査は別の操作として記録する。

bundleのルートで次を実行すると、入力を一時コピーして保存証明書を監査する。
\begin{verbatim}
python3 -B verification/run_v14_checks.py
\end{verbatim}
指標の全再計算、選択テンソル縮約の独立な再計算、またはTateの新規再構成には、
同じスクリプトの \texttt{--help} にある追加オプションを用いる。
検算にはPython 3.12.14とNumPy 2.3.5を用いた。
計算は整数・有理数・明示された有限体で行い、浮動小数点の近似階数判定を使わない。

原稿の幾何学的な同定と既存定理は引き続き証明の入力である。
有限計算の成功は、原稿全体の全ての幾何学的・圏論的議論を独立に再認証したことを意味しない。
また、文献の新規性や優先権の全面調査は本版の作業範囲に含めていない。


\section{第15版の証明書・保存範囲・再現手順}
\label{app:v15-reproducibility}

第15版の追加研究は継続用bundle v3の \texttt{source\_v15/} に収めた。
添付由来の \texttt{source\_v13/} と、前回の \texttt{source\_v14/} は保存し、
旧原稿も比較と再構成の基準として保持した。
第15版の日本語TeXは単体でLuaLaTeXにより組版でき、
基準となる第14版と追加章からの再組立てにも対応する。

今回の証明の入力は次のように区別する。
指標は $F$ 次数15まで全重みをEuler積とNewton--Adamsで新規に照合した。
接触部分空間の164方向とKoszul核の359方向は、整係数テンソル公式、
公開整数入力、法1009の非零小行列式により特性零の独立性を証明した。
採用した全縮約は別の辺番号付けでも再計算している。
三点の164接触列と二つの補元、14個と1個の二階補正は有理数で再構成した。
最後に、三点のTate係数も保存値の補間によらず初項から再帰で求めた。

\nolinkurl{verification/run_v15_checks.py} は入力を一時コピーして監査し、
JSONとテキストの結果を作る。
\begin{verbatim}
python3 -B verification/run_v15_checks.py
\end{verbatim}
この既定実行は保存証明書の監査を中心とする。
全重み、テンソル縮約、またはTate再帰そのものの新規再計算は、
同スクリプトの \texttt{--help} に表示される追加オプションで区別して指定する。
Pythonのassertionを無効化する最適化オプションは使わない。
依存はPython標準ライブラリとNumPyであり、
検算時の版はPython 3.12.14、NumPy 2.3.5である。

高速監査の成功を、再帰を初めから作り直した結果と同一視しない。
また、第15版は通常20次商の次元を1または2という範囲に留めており、
その完全同定を実際の20次消去の前提にしていない。
幾何学的同定と既存定理は引き続き数学的入力であり、有限計算の成功が
既存原稿の全ての圏論的・幾何学的議論の独立な再認証を意味するわけではない。
文献の新規性・優先権についての全面調査も、今回の計算研究の主張には含めていない。


\section{第16版の証明依存・再現手順・継続点}
\label{app:v16-reproducibility}

第16版の追加研究は継続用bundle v4の \texttt{source\_v16/} に収めた。
第13版、第14版、第15版の研究ソースはバイト単位で保存し、
各版のTeX、PDF、証明書と前版の索引も比較可能な形で保持した。
第16版のTeXは単体でLuaLaTeXにより組版でき、
変更していない第15版と追加章からの再組立ても検査した。

\paragraph{構造証明.}
Pl\"ucker素イデアル、普遍方向の錐上消失、adjugate恒等式、
単元制限環、横断三次項による単射、普遍勾配のheightを順に扱った。
三元三次不変式環の生成定理とCohen--Macaulay環の正則列判定は
明示した一次資料を用いる。
新規性・優先権の全面的調査を実施したという主張は含めない。
記号恒等式の検算例と、証明そのものの論理的入力は本文で区別した。

\paragraph{有限計算の証明書.}
二十二次の指標は、$F$ 次数17、$a$ 次数22までの全重みを
Euler積とNewton--Adams再帰で照合した。
通常像の429方向は公開整数入力上の非零minorで独立性を証明し、
選択した全 $429^2$ 成分を別の縮約実装でも再構成した。
全429方向の二点値、143個の積方向の因数分解、76個の二階補正を
有理数で検算した。
実際のTate値は四入力五回の新規Maurer--Cartan再帰による。
保存ブロックからの行列式再計算と、再帰そのものの再実行は区別する。

\paragraph{再現の入口.}
\begin{verbatim}
python3 -B verification/verify_integrity.py
python3 -B verification/run_v16_checks.py --help
python3 -B verification/run_v16_checks.py
\end{verbatim}
既定の監査と、全指標・全縮約・Tate再帰の再計算はオプションで分けた。
既定の成功が全ての重い再帰をその場で再実行したことを意味しない。
assertionを有効にして実行し、Pythonの最適化オプションは使わない。
入力を一時コピーして検査するため、配布する確定済みソースは変更されない。
依存するPython、NumPy、および日本語LuaLaTeXの条件は同梱の手順に示した。

\paragraph{次回に引き継ぐ境界.}
現在確定しているのは24次開始の偶正規形、20次商2次元、22次商1次元、
二項下位jetに対する全次数の $I_N=J_N$、および自由部分加群の単射である。
全商の同定、24次以降の実際の消滅、残る第2Koszul微分の階数は未決定である。
研究状態、主張台帳、本文の命題・ページ索引、コード索引、
次回の具体的な手順をbundleに収めた。


\section{第17版の全次数構造と二十四次計算の再現範囲}
\label{app:v17-reproducibility}

第17版では、通常商の構造を証明する四つの章と、二十四次を確定する六つの章を加えた。
継続用bundleは第5版であり、原稿の第17版とは版番号を区別する。
\nolinkurl{source_v13/}、\nolinkurl{source_v14/}、
\nolinkurl{source_v15/}、\nolinkurl{source_v16/} は不変の基準として保持した。
今回の新しい脚本、入力、結果、証明原稿は \nolinkurl{source_v17/} に置く。

\subsection{構造理論の証明依存}

行列束とadjugateの恒等式は、対称行列15変数と交代行列10変数を独立変数とした
整数係数多項式の恒等式として検証する。
三元三次式の全共変量加群についてはFisherの
\cite[Theorem~2.7(i)]{FisherHessian2012}を用い、判別式を可逆にしない全係数環上で
normal射影を構成する。
Heisenberg作用のcommutantとHesse族の優勢性は、階数8と行列式118098で確認する。
頂点支持の証明は、\v{C}ech--Koszul双複体の二つの収束するスペクトル系列を比較する。
有限次元性はNoether有限生成性と指定二重次数から従う。
これらの数学的推論を、有限標本の行列の階数から導いたとは扱わない。

\subsection{二十四次の証明書}

全指標計算は $F$ 次数19、$a$ 次数24までの重みをEuler積とNewton再帰で照合し、
Schur--Weyl型の交代和と独立なPieri列挙から各次元を求める。
接触像では1120個の有理点を使い、1365次元の $PU_{18}$ を消した上で
1091本の通常接触方向の非零minorを構成する。
別実装は905個の相異なる因子を新規に縮約し、選択小行列の1190281成分を全て照合する。
保存された因子から全行列を再構成する軽量監査も用意した。
軽量監査と新規テンソル縮約は、実行記録上も区別する。

通常商の上界にはnormal sliceの独立性を使う。
有限体で零になった値、探索の停滞、KoszulのEuler和だけを上界に用いない。
二階補正の所属は全次数のnormal射影零と
$\mathcal K_{24}^{\mathrm{nor}}=0$ を組み合わせて証明する。
四入力での310本の補正の有限評価は追加の照合であり、
そのうち二本が点ごとには非零であることも記録した。

\subsection{Tate計算の実施範囲}

四つの完全な点入力で五回の新規再帰を行った。
このうち一入力は、共通分母整数演算とは別の\texttt{Fraction}演算でも再帰し、
全20定数ブロックと24次以下の全25行列式係数を比較した。
さらに二つの入力について、$\mathbb Q[\varepsilon]/(\varepsilon^4)$ 上で
法三次jetを最初から生成した。
この二回は、点での値から補間した計算ではない。
法三次再帰の保存値監査は独立な\texttt{Fraction}実装で全200係数を再構成し、
入力テンソル・交代行列・実行脚本のhashも確認する。

有理数体上のMaurer--Cartan検査、有限体上の独立な32次元外積作用素検査、
二つの行列式算法による照合を区別して記録する。
過去の切断済み $A_{15},A_{16}$ を全項の再帰として再利用していない。
二十六次へ進む場合は、全 $A_1,\ldots,A_{18}$ と、
曲率 $R_{19},\ldots,R_{22}$ の内部$x$次数 $4,3,2,1$ 以下が新たに必要になる。

\subsection{同梱された入口}

bundleを展開したディレクトリで次を実行する。
Pythonのassertを有効に保ち、\texttt{-O}、\texttt{-OO}を用いない。
\begin{verbatim}
python3 -B verification/verify_integrity.py
python3 -B verification/run_v17_checks.py
python3 -B verification/run_v17_checks.py --fresh-normal-jets
\end{verbatim}
既定の10項目は、指標保存監査、KoszulのPieri計算、通常像の保存監査、
構造恒等式、24次台帳、G12基底の新規厳密計算、normal検出子、
下位補正の保存監査、完全点Tate保存監査、法三次Tate保存監査である。
入口は五つのソース木を一時コピーし、実行前後の元ファイルのhashを比較する。
\begin{verbatim}
python3 -B verification/run_v17_checks.py --full-characters
python3 -B verification/run_v17_checks.py --recompute-graphs
python3 -B verification/run_v17_checks.py --recompute-detectors
\end{verbatim}
これらのオプションは、それぞれ全重み、接触minorの独立新規縮約、
四入力の310補正の新規計算を追加する。
完全点のTate再帰を最初から実行するコマンドは
\nolinkurl{source_v17/README_ja.txt} に記す。

\subsection{到達点と残る課題}

通常商は自由な $\Theta_8$ 成分と有限次元の頂点部分に分裂した。
しかし有限補正多項式 $E_{\mathrm{vert}}$ の全係数と最高次数、
自由部分と一致し始める閾値は未計算である。
二十四次の確定から二十六次の頂点部分の消失を外挿しない。
二十六次の既知の自由方向は
$[\Theta_8U_4^3J_6]$、$[\Theta_8J_6^3]$ である。
実際の係数については、raw Tate項、正規化後の項、normal射影、頂点部分を区別して調べる。
全次数でのHessian正規形と、残る $H_1,H_2$ の個別次元は未決定である。
本版は、新規性と優先権についての網羅的文献調査を行ったとは主張しない。


\section{第18版の証明依存と再現範囲}
\label{app:v18-reproducibility}

本稿の第18版に対応する継続用bundleは第6版である。
第13版から第17版までのソース木、原稿、証明書と索引は基準資料として保持し、
今回の追加を \nolinkurl{source_v18/} にまとめる。
過去の節の「未決定」という記述は当該版の到達点を表し、
現在の到達点は冒頭、今回の七つの節、および今後の課題に記載する。

\subsection{新しい証明の入力}

全次数normal消滅の証明は、既存の固定平面のPicard直線とnormal coneの幾何を入力とする。
その上でHesse pencilの形式一意性を証明し、全多項式共変量へ延長する。
新しい多項式監査はHessian、Hesseパラメータの変更式、$j$ の有理式とその微分、
および優勢性を示す十行十列行列を確認する。
既存の幾何学的証明全体を有限コードで再認証したという意味ではない。

接触変換則はSchur chart、横断次数の重み、三元End共変量のスカラー性から導く。
有限次数上界にはDub\'eの斉次版定理\cite[Theorem~8.2]{Dube1990}を用いる。
45変数全体を保ち、$F$ を係数体に移していない。
Hessianのsocle恒等式はEuler式、adjugate、$aq=0$ の三つの普遍恒等式から導く。
その微分の頂点支持には第17版のSchur chart上のČech--Koszul次数評価を使う。
これらの記述証明と、有限多項式の監査記録は別の役割を持つ。

\subsection{26次計算の実施範囲}

全指標は $F$ 次数21、$a$ 次数26までのEuler積とNewton--Adams再帰を照合する。
Schur重複度、Pieri列挙、全Weyl次元和、および旧版との重複を確認する。
通常26次商を決める完全な接触minorは今回構成していない。
96列の部分minorの探索とその独立監査は、完全な階数証明と区別する。

法三次Tate計算は二入力について最初から行い、
全 $A_1,\ldots,A_{18}$ と必要な切断 $A_{19},\ldots,A_{22}$ を再構成する。
旧版で切断した $A_{17},A_{18}$ を流用していない。
独立な保存監査は二つの行列式算法を使い、二入力の全216係数を再構成する。
通常の点における26次Tate係数の完全再帰は今回の実施範囲に含めない。
準備した完全点ドライバは未実行の補助コードとして明記する。

\subsection{再開用の入口と文書上の訂正}

bundleを展開したディレクトリで次を実行する。
Pythonのassertを有効に保ち、\texttt{-O}、\texttt{-OO}を用いない。
\begin{verbatim}
python3 -B verification/verify_integrity.py
python3 -B verification/run_v18_checks.py
python3 -B verification/run_v18_checks.py --fresh-normal-jets
python3 -B verification/run_v18_checks.py --full-characters
\end{verbatim}
検証入口は六つのソース木を一時コピーし、元のファイルのhashが変わらないことを確認する。
既定の保存監査、新規法三次再帰、全重みの再計算は、実行記録上も区別する。
正確な入力、出力、依存関係と所要時間は
\nolinkurl{source_v18/README_ja.txt} と再現報告にまとめた。

第17版の \nolinkurl{source_v17/README_ja.txt} にある完全点Tateの新規再帰例は、
\texttt{--audit} を付ける点に文書上の誤りがあった。
当該ドライバの \texttt{--audit} は既存JSONの保存ブロック監査を選ぶため、
新規再帰ではこのフラグを外す。
保存されている過去の新規計算記録そのものは変更せず、
現在のREADMEに正しい二種類のコマンドを示す。

\subsection{未解決事項を引き継ぐ}

全次数normal消滅と、normal消滅部分での明示的有限決定性は今回の新しい結論である。
$E_{\mathrm{vert}}=t^{10}$、$\mathcal K_{26}^{\mathrm{nor}}=0$、
実際の26次項の通常像への所属、無限形式変換の収束・代数化は残る。
従って26次開始の既存正規形を28次開始へ更新していない。
新規性や優先権について網羅的な文献調査を行ったとは主張しない。


